%% kapa: exam=1 round=25 year=2014 period=2 form=A subject=accounting subject_ko=회계학 seq=21 number=61
%% source: Q-Net 감정평가사 기출문제 (cst00309) artlSeq=5101207 file=extracted/1008/5101207/2014년 제25회 감정평가사 1차시험 기출문제/2교시(A형).hwp
%% conversion: kordoc 4.24.51 -> tools/gen1.py (auto); review=proofread-sonnet; flags=none
\begin{question}{61}
(주)서울은 20×1년 1월 1일에 ￦10,000에 기계장비를 취득하였다. 이 기계장비의 내용연수는 10년이며, 잔존가치는 없으며, 정액법을 이용하여 감가상각하였다. 20×1년 말과 20×2년 말에 (주)서울은 이 기계장치에 대한 손상 징후가 있다고 판단하여 손상검사를 실시하였는데, 이에 대한 정보는 아래와 같다.
\begin{framedtext}
\begin{tabular}{@{}l@{\hspace{3em}}r@{\hspace{2em}}r@{}}
 & \uline{20×1.12.31} & \uline{20×2.12.31} \\
순공정가치 & ￦7,200 & ￦7,750 \\
사용가치 & 8,100 & 7,600
\end{tabular}
\end{framedtext}
이 기간 중 내용연수, 잔존가치 및 감가상각방법에 변화가 없었다면, 20×2년 포괄손익계산서에 인식할 내용은? \par
\begin{choices}
  \item 손상차손 ￦350
  \item 손상차손 ￦400
  \item 손상차손 ￦500
  \item 손상차손환입 ￦550
  \item 손상차손환입 ￦1,200
\end{choices}
\end{question}
\answer{4}
